% ----------------------------------------------------------
\chapter{Código do Classificador Operativo de Gargalos}
\label{ap:regras_classificador}
% ----------------------------------------------------------

Este apêndice apresenta a rotina de classificação baseada em limiares físicos e lógicos desenvolvida em Python. O algoritmo avalia as métricas de rádio de cada registro de conexão telemétrica e categoriza a amostra operacional em perfis mutuamente exclusivos de gargalo físico ou ideal.

\begin{lstlisting}[language=Python, caption={Rotina de Classificação de Gargalos de Rede}, label={lst:classificador_gargalos}]
def classify_bottlenecks(df):
    """
    Classifica cada registro de telemetria nos perfis de gargalo
    de rede ou operacao ideal baseados nos limiares heuristicos.
    """
    # 1. Gargalo G1: Baixa Qualidade de Enlace (sinal < -75 dBm)
    df['gargalo_sinal'] = df['signal'] < -75
    
    # 2. Gargalo G2: Congestionamento do Meio (clientes ativos >= 15)
    df['gargalo_congestionamento'] = df['clientes_ativos_por_ap_minuto'] >= 15
    
    # 3. Gargalo G3: Interferencia / SNR Ruim (ruido > -90 dBm ou SNR < 20 dB com sinal saudavel)
    df['gargalo_interferencia'] = (df['noise'] > -90) | ((df['signal'] >= -75) & (df['SNR'] < 20))
    
    # 4. Gargalo G4: Instabilidade de Roaming (mais de 5 transicoes no dia)
    df['gargalo_roaming'] = df['total_roaming'] > 5
    
    # 5. Gargalo G5: Gargalo Cabeado (throughput do AP >= 90 Mbps - limite Fast Ethernet)
    df['gargalo_cabo'] = df['tp_agregado_ap'] >= 90
    
    # Estado Ideal: Amostras sem ocorrencia de nenhum dos limitantes anteriores
    df['sem_gargalo'] = ~(
        df['gargalo_sinal'] | 
        df['gargalo_congestionamento'] | 
        df['gargalo_interferencia'] | 
        df['gargalo_roaming'] | 
        df['gargalo_cabo']
    )
    
    return df
\end{lstlisting}
